% Apêndices (pós-textuais, depois das referências). Cada \chapter vira
% "APÊNDICE X – TÍTULO" em página própria; o conteúdo começa na página seguinte.

\chapter{Mapa de Implementação}
\label{ap:mapa}

Este apêndice relaciona os componentes descritos no Capítulo~\ref{cha:arquitetura} aos arquivos do repositório e reúne os detalhes de implementação que não sustentam a argumentação do corpo do texto.

\section{Tecnologias da implementação}

A Tabela~\ref{tab:stack} lista as tecnologias empregadas. O sistema usa Python como linguagem base e não depende de framework web: o servidor HTTP é da biblioteca padrão, com tratamento concorrente de requisições (\texttt{socketserver.ThreadingMixIn}). A simulação histórica clássica é executada por um motor de \textit{backtesting} próprio, independente de bibliotecas externas.

\begin{table}[htbp]
\centering
\caption{Tecnologias da implementação do Hedge-fund-lab.}
\label{tab:stack}
\small
\begin{tabular}{>{\raggedright\arraybackslash}p{3.8cm} >{\raggedright\arraybackslash}p{11.2cm}}
\toprule
\textbf{Camada} & \textbf{Tecnologia utilizada} \\
\midrule
Servidor HTTP & \texttt{http.server}/\texttt{socketserver} (biblioteca padrão), \texttt{dashboard/server.py} \\
Frontend & HTML/CSS/JS + Chart.js (\texttt{dashboard/}) \\
Persistência de séries & SQLAlchemy, URL configurável, padrão SQLite (\texttt{src/db}) \\
Livro das carteiras & SQLite direto, eventos encadeados (\texttt{dashboard/manual\_portfolio\_api.py}) \\
Orquestração (ativo único) & LangGraph (\texttt{src/agents}) \\
Carteira multiativos & \texttt{scripts/multiasset\_radar.py}, \texttt{dashboard/wallet\_manager\_api.py} \\
Replay contínuo & \texttt{scripts/continuous\_replay*.py}, \texttt{scripts/export\_continuous\_replay*.py} \\
Contratos de resposta & Pydantic (esquemas estruturados das saídas dos modelos) \\
Modelo de linguagem & \texttt{gemini-3.8-flash} (Google), via transporte HTTP próprio \\
\textit{Backtesting} & motor próprio (\texttt{src/backtesting}) \\
Estratégias & \texttt{src/strategies}, \texttt{dashboard/rule\_portfolios.py} \\
Indicadores & \texttt{src/indicators} \\
Pipeline de dados & \texttt{src/pipeline} (incl. \texttt{b3\_official.py}) \\
Experimentos & \texttt{src/experiments} \\
Testes & pytest (Poetry) e verificações por \texttt{assert} em \texttt{tests/} \\
\bottomrule
\end{tabular}
\\[2pt]
{\footnotesize Fonte: elaborado pelos autores, a partir do repositório do projeto.}
\end{table}

\section{Componentes e arquivos}

A Tabela~\ref{tab:mapa-arquivos} associa cada componente do Capítulo~\ref{cha:arquitetura} aos arquivos que o implementam. O caminho multiativos é implementado em \texttt{scripts/} e \texttt{dashboard/}, sem alterar os módulos de \texttt{src/} que integram a proveniência do protocolo de ativo único. As versões v2 e v3 e o gestor consultivo reutilizam módulos anteriores por vinculação de módulo (\texttt{pilot\_paper.bind}), sem editar as fontes seladas.

\begin{table}[htbp]
\centering
\caption{Mapeamento entre componentes e arquivos do repositório.}
\label{tab:mapa-arquivos}
\small
\begin{tabular}{>{\raggedright\arraybackslash}p{4.0cm} >{\raggedright\arraybackslash}p{11.0cm}}
\toprule
\textbf{Componente} & \textbf{Arquivos} \\
\midrule
Dados oficiais e ajuste & \texttt{src/pipeline/b3\_official.py}, \texttt{extract.py}, \texttt{transform.py}, \texttt{load.py}, \texttt{snapshot.py}, \texttt{flows.py} \\
Calendário & \texttt{src/backtesting/b3\_calendar.py} (\texttt{B3Calendar}) \\
Base de mercado e \textit{bundle} & \texttt{dashboard/market\_base\_api.py}, \texttt{scripts/continuous\_replay\_data.py} \\
Livro contábil & \texttt{dashboard/manual\_portfolio\_api.py}, \texttt{scripts/multiasset\_portfolio.py} (núcleo contábil) \\
Comitê e risco por ativo & \texttt{scripts/multiasset\_radar.py}, \texttt{src/agents/technical\_analyst.py}, \texttt{technical\_evidence.py}, \texttt{features.py}, \texttt{feature\_semantics.py} \\
Gestor e compilador & \texttt{dashboard/wallet\_manager\_api.py}, \texttt{scripts/run\_global\_pm\_v2.py} (transporte e diário reutilizados) \\
Planejador de execução & \texttt{dashboard/auction\_orders.py} \\
Proventos & \texttt{dashboard/dividends\_api.py}, \texttt{scripts/continuous\_replay\_dividends.py} \\
Relógio causal & \texttt{scripts/replay\_multiasset.py}, \texttt{tests/check\_multiasset\_replay.py} \\
Carteiras de referência & \texttt{dashboard/rule\_portfolios.py}, \texttt{src/strategies/sma\_cross.py}, \texttt{bollinger\_bands.py}, \texttt{src/backtesting/portfolio.py} \\
Executores do \textit{replay} & \texttt{scripts/continuous\_replay.py}, \texttt{continuous\_replay\_v2.py}, \texttt{continuous\_replay\_v3.py}, \texttt{continuous\_replay\_runtime.py} \\
Auditorias de paridade & \texttt{scripts/audit\_replay\_v2\_comparison.py}, \texttt{audit\_replay\_v3\_comparison.py} \\
Qualificação sem rede & \texttt{tests/check\_continuous\_replay\_v2.py}, \texttt{check\_continuous\_replay\_v3.py} \\
Exportação e publicação & \texttt{scripts/export\_continuous\_replay*.py}, \texttt{build\_replay\_explorer.py}, \texttt{publish\_replay.py} \\
Ciclo diário & \texttt{dashboard/daily\_cycle.py} \\
Servidor e interface & \texttt{dashboard/server.py}, \texttt{dashboard/} (HTML, CSS e JS) \\
Implantação & \texttt{Dockerfile}, \texttt{railway.json} \\
\bottomrule
\end{tabular}
\\[2pt]
{\footnotesize Fonte: elaborado pelos autores, a partir do repositório do projeto.}
\end{table}

\section{Detalhes de persistência}

As séries do pipeline são persistidas pelo módulo \texttt{src/db}, composto pela configuração de conexão (\texttt{connection.py}) e pela definição do esquema relacional (\texttt{models.py}); o endereço do banco é configurável (\texttt{src/config.py}), com SQLite como valor padrão.

O livro contábil das carteiras (identidade \texttt{HFL-MANUAL-PORTFOLIO-V1}) grava cada evento em uma tabela SQLite com gatilhos que recusam \texttt{UPDATE} e \texttt{DELETE}. As escritas usam \texttt{BEGIN IMMEDIATE} e \texttt{synchronous=FULL}. Valores monetários e quantidades são tratados com \texttt{Decimal} e arredondados a centavos pelo núcleo contábil reutilizado de \texttt{scripts/multiasset\_portfolio.py}.

\section{Códigos de erro do relógio causal}

Os controles de causalidade descritos na Seção~\ref{sec:arq-relogio} interrompem a execução com os seguintes códigos:

\begin{itemize}
  \item \texttt{decision\_input\_lookahead}: a verificação preliminar das versões v2 e v3, executada sob bloqueio de rede (\texttt{no\_network()}), encontrou uma sessão de entrada igual ou posterior ao dia-alvo;
  \item \texttt{input\_changed\_before\_transport}: a entrada de um ativo, relida imediatamente antes da requisição HTTP, difere da versão registrada;
  \item \texttt{manager\_context\_changed\_before\_transport}: o mesmo para o contexto do gestor;
  \item \texttt{campaign\_identity\_changed}: já existe, na raiz da campanha, um manifesto diferente do que seria gravado (Seção~\ref{sec:arq-replay}).
\end{itemize}

\section{Reservas de orçamento e tarifas}

A reserva gravada antes de cada tentativa HTTP é de US\$~0,110592 por chamada de analista ou de risco e de US\$~0,159744 por chamada do gestor, valores calculados pelos limites de \textit{bytes} de entrada e \textit{tokens} de saída às tarifas mais altas anunciadas pelo provedor. Ao receber a resposta, o custo é estimado com US\$~0,75 por milhão de \textit{tokens} de entrada e US\$~3,75 por milhão de \textit{tokens} de saída e de raciocínio; sem uso válido, a reserva integral permanece contabilizada. Esperas de reenvio superiores a 30 segundos interrompem a execução. O controle é implementado em \texttt{scripts/continuous\_replay\_runtime.py}, é independente da estratégia e a auditoria correspondente (\texttt{audit\_budget}) está nos exportadores.

\section{Constantes da versão v3}

A Tabela~\ref{tab:constantes-v3} reúne as constantes determinísticas de \texttt{scripts/continuous\_replay\_v3.py} referidas na Seção~\ref{sec:arq-evolucao}.

\begin{table}[htbp]
\centering
\caption{Constantes determinísticas da versão v3.}
\label{tab:constantes-v3}
\small
\begin{tabular}{l r p{6.4cm}}
\toprule
\textbf{Constante} & \textbf{Valor} & \textbf{Função} \\
\midrule
\texttt{MAX\_ASSET\_WEIGHT} & 0,30 & Teto de peso projetado em compra por COMPRA \\
\texttt{DIP\_REBALANCE\_MAX\_WEIGHT} & 0,22 & Teto de peso em compra de recomposição em MANTER \\
\texttt{MIN\_TRADE\_BRL} & 2.000,00 & Valor bruto mínimo de compra (R\$) \\
\texttt{CORE\_FLOOR\_WEIGHT} & 0,08 & Peso mínimo remanescente após venda técnica normal \\
\texttt{MAX\_NORMAL\_SELL\_FRACTION} & 0,45 & Fração máxima da posição vendida por pregão \\
\texttt{ANTI\_WHIPSAW\_COOLDOWN\_SESSIONS} & 8 & Janela contra recompra mais cara (pregões) \\
\texttt{CASH\_DRAG\_THRESHOLD} & 0,25 & Caixa/patrimônio que habilita a recomposição \\
\texttt{STRENGTH\_TRIM\_MAX\_FRACTION} & 0,35 & Fração máxima de uma realização em força \\
\texttt{STRENGTH\_TRIM\_MIN\_WEIGHT\_AFTER} & 0,12 & Peso mínimo após realização em força \\
\texttt{TRIM\_MIN\_WEIGHT\_BEFORE}/\texttt{\_AFTER} & 0,25/0,15 & Limites do \textit{Trim-to-Fund}, herdados da v2 \\
\bottomrule
\end{tabular}
\end{table}

\section{Regras detalhadas das versões v2 e v3}

Esta seção detalha as regras resumidas na Seção~\ref{sec:arq-evolucao}.

\subsection{Versão v2}

A v2 acrescenta quatro intervenções, todas impostas tanto no \textit{prompt} do gestor (\texttt{PROMPT\_V2}) quanto no compilador (\texttt{compile\_proposal\_v2}):

\begin{itemize}
  \item \textbf{\textit{Trim-to-Fund}:} venda parcial de ativo com recomendação MANTER ou COMPRA, admitida quando o peso antes da operação supera 25\% do patrimônio, o peso projetado permanece em ao menos 15\% e a cesta contém uma compra válida ou o ativo já excede 30\%.
  \item \textbf{Teto de concentração e ordem mínima:} nenhuma compra pode levar o peso projetado do ativo acima de 30\% do patrimônio anterior à operação (\texttt{MAX\_ASSET\_WEIGHT}), nem ter valor bruto inferior a R\$~2.000,00 (\texttt{MIN\_TRADE\_BRL}); vendas não têm valor mínimo.
  \item \textbf{Veto assimétrico (R7):} o veredito \texttt{VETADO} passa a bloquear apenas compras; se o consenso técnico for VENDA, a recomendação permanece VENDA sob veto, permitindo reduzir ou zerar a posição.
  \item \textbf{Lentes analíticas:} cada analista recebe o \textit{prompt} técnico v5 acrescido de uma lente --- tendência estrutural (analistas 1 e 2), reversão por Bollinger e RSI (3 e 4) e regime de volatilidade como árbitro (5).
\end{itemize}

Além disso, o compilador deixa de invalidar o pregão diante de uma violação de regra: em modo não estrito, reduz a quantidade ou converte a ação em \texttt{HOLD} e registra cada ajuste com \textit{ticker}, lado, quantidades de origem e destino e o motivo (por exemplo, \texttt{MAX\_ASSET\_WEIGHT\_CAP}, \texttt{MIN\_TRADE\_BRL}, \texttt{TRIM\_WEIGHT\_AFTER\_FLOOR}). Violações estruturais --- universo incompleto ou com \textit{tickers} repetidos e quantidades inválidas --- continuam invalidando o pregão.

\subsection{Versão v3}

A v3 mantém as regras da v2 e acrescenta quatro grupos de mecanismos, parametrizados pelas constantes da Tabela~\ref{tab:constantes-v3}.

\begin{itemize}
  \item \textbf{Cinco lentes especialistas e travas de extremo.} Cada analista recebe uma lente única: Bandas de Bollinger, exaustão por RSI e MACD, valor relativo à média de 50 pregões, estrutura de longo prazo pela média de 200 pregões e assimetria entre o espaço até as bandas superior e inferior. Cada lente explicita faixas em que determinado sinal é proibido. Duas travas atuam apenas sobre consensos de exatamente três votos: \texttt{ANTI\_BOTTOM\_SELL\_GUARD} reverte VENDA para MANTER quando o RSI é inferior a 40 ou \texttt{bb\_lower\_gap} é inferior a 0,035, desde que \texttt{sma200\_gap} seja superior a $-0{,}08$; \texttt{ANTI\_TOP\_CHASE\_GUARD} reverte COMPRA para MANTER quando o RSI é superior a 62 ou \texttt{bb\_upper\_gap} é superior a $-0{,}02$.
  \item \textbf{Memória de posição, piso núcleo, venda escalonada e \textit{anti-whipsaw}.} O contexto do gestor passa a incluir, por ativo, a quantidade, o preço médio de compra, o retorno não realizado, a sessão e o preço da última venda e os pregões decorridos desde ela, reconstruídos do livro. Em venda técnica normal (VENDA com risco aprovado), o compilador limita a quantidade a 45\% da posição, arredondada para cima, e preserva peso remanescente de ao menos 8\% (ajuste \texttt{CORE\_FLOOR\_STEP\_SELL}); a saída integral fica restrita à venda sob veto (R7). É recusada a compra de ativo vendido nos últimos oito pregões a preço superior ao da última venda (\texttt{ANTI\_WHIPSAW\_HIGHER\_REBUY}).
  \item \textbf{\textit{Trim-on-Strength} e \textit{Trim-to-Fund}.} Além do \textit{Trim-to-Fund} da v2, admite-se a realização parcial em força de ativo em MANTER ou COMPRA quando ele está esticado (RSI $\geq$ 63 ou \texttt{bb\_upper\_gap} $\geq -0{,}02$), em lucro de ao menos 2\% sobre o preço médio e com peso acima de 14\%, limitada a 35\% da posição e a peso remanescente de 12\%, sem exigir compra na cesta.
  \item \textbf{Recomposição contra o arrasto de caixa (\textit{Dip-Rebalance}).} Quando o caixa do início do pregão supera 25\% do patrimônio, admite-se compra de ativo em MANTER com risco aprovado e em desconto (\texttt{sma50\_gap} negativo com RSI de até 46, ou \texttt{bb\_lower\_gap} de até 0,04), sujeita ao teto de 22\%, à ordem mínima e à regra \textit{anti-whipsaw}.
\end{itemize}



\chapter{Caminho de Ativo Único e Camada de Experimentos}
\label{ap:ativo-unico}

Este apêndice descreve o caminho de decisão de ativo único, construído para o protocolo confirmatório de ativo único (PETR4), planejado na etapa anterior para avaliar H2 (Seção~\ref{sec:met-protocolo}), e a camada de experimentos que o sustenta. Suas fases de desenvolvimento foram executadas, mas o Teste Final não; por isso o caminho não aparece nos resultados.

\section{Dados do protocolo de ativo único}

A série de PETR4 do protocolo vem de um \textit{snapshot} oficial do COTAHIST, ajustado pelo fator do \texttt{yfinance}, com 2.649 barras, uma por sessão oficial de 04/01/2016 a 31/08/2026. Dela, 1.519 sessões (de 12/01/2018 a 28/02/2024), cada uma com ao menos 504 sessões de histórico, formam o conjunto elegível de âncoras; quatro dessas sessões, escolhidas por quantis e sem olhar o mercado, compõem o conjunto de diagnóstico \texttt{H\_real}. O calendário de negociação (\texttt{B3Calendar}, em \path{src/backtesting/b3_calendar.py}) aplica as regras históricas de feriados da B3 e reproduz as 2.649 sessões oficiais.

\section{Indicadores, estratégias e motor de \textit{backtesting}}

Os indicadores técnicos são implementados em \texttt{src/indicators}: médias móveis simples (\texttt{sma.py}), Bandas de Bollinger (\texttt{bollinger.py}), Índice de Força Relativa (\texttt{rsi.py}, formulação canônica de Wilder) e MACD (\texttt{macd.py}).

As estratégias residem em \texttt{src/strategies}, com uma interface comum em \texttt{base.py}: \textit{Buy and Hold} (\texttt{buy\_and\_hold.py}), carteira igualmente ponderada (\texttt{equal\_weight.py}), cruzamento de médias móveis (\texttt{sma\_cross.py}), Bandas de Bollinger (\texttt{bollinger\_bands.py}), mínima variância (\texttt{min\_variance.py}) e o controle clássico da família de indicadores (\texttt{indicator\_family\_control.py}). Esse último combina SMA, Bandas de Bollinger, RSI e MACD em um escore discreto no conjunto $\{-1, 0, +1\}$, servindo de \textit{benchmark} clássico para a comparação com o sistema multiagente de ativo único. No caminho multiativos, \texttt{dashboard/rule\_portfolios.py} aplica \texttt{sma\_cross.py} (50/200) e \texttt{bollinger\_bands.py} (20 períodos, 2 desvios) a parcelas independentes de 20\% do capital por ativo, e as carteiras de pesos iguais e de mínima variância utilizam as classes \texttt{EqualWeightPortfolio} e \texttt{MinVariancePortfolio} de \texttt{src/backtesting/portfolio.py}.

O motor de \textit{backtesting} (\texttt{src/backtesting}) centraliza a simulação de carteiras, o cálculo de métricas e a incorporação de custos: \texttt{engine.py} conduz a simulação, \texttt{portfolio.py} modela a carteira, \texttt{costs.py} aplica custos de transação e \texttt{metrics.py} concentra as métricas de risco e retorno. O calendário (\texttt{b3\_calendar.py}) alinha as decisões às sessões oficiais. A integração com o sistema multiagente de ativo único é feita por \texttt{agent\_engine.py} e \texttt{daily\_agent.py}, que conectam as decisões dos agentes à execução simulada, e por \texttt{arena.py}, que organiza a comparação entre estratégias. A métrica primária de avaliação do sistema multiagente (\textit{Sharpe Ratio}) é definida de forma congelada em \texttt{metrics.py}, sob o identificador \texttt{H2\_SCIENTIFIC\_SHARPE\_DEFINITION\_V1}, cuja especificação é detalhada no Capítulo~\ref{cha:metodologia}. O Replay Contínuo multiativos não utiliza \texttt{engine.py}.

\section{Sistema multiagente de ativo único}
\label{sec:arq-ativo-unico}

A lógica de decisão de ativo único é modelada como um sistema multiagente coordenado pelo LangGraph, no módulo \texttt{src/agents}. O grafo é linear e acíclico (Figura~\ref{fig:grafo-langgraph}): os agentes são executados em sequência, sem debate cíclico nem reavaliação iterativa. Cada nó corresponde a um papel especializado, e as arestas definem a ordem de execução. A aresta condicional após o gestor de risco só encaminha a decisão ao gestor de portfólio quando o veredito consolidado é \texttt{APROVADO}.

\begin{figure}[htbp]
\centering
\begin{tikzpicture}[
  node distance=5mm,
  box/.style={draw, rounded corners=2pt, minimum height=7mm, align=center, inner sep=3pt, font=\small},
  term/.style={draw, circle, minimum size=7mm, align=center, font=\small\bfseries, inner sep=1pt},
  arr/.style={-Stealth, thick},
]
\node[term] (start) {início};
\node[box, below=of start] (ta) {\texttt{technical\_analyst}\\análise técnica};
\node[box, below=of ta] (rm) {\texttt{risk\_manager}\\gestão de risco};
\node[box, below left=10mm and 28mm of rm] (pm) {\texttt{portfolio\_manager}\\gestão de portfólio};
\node[term, below=of rm] (end) {fim};
\draw[arr] (start) -- (ta);
\draw[arr] (ta) -- (rm);
\draw[arr] (rm) -- node[above, sloped, pos=0.55, font=\footnotesize] {veredito = \texttt{APROVADO}} (pm);
\draw[arr] (pm) -- (end);
\draw[arr] (rm) -- node[right, font=\footnotesize] {caso contrário} (end);
\end{tikzpicture}
\caption{Grafo de decisão do caminho de ativo único (LangGraph): sequência linear \texttt{technical\_analyst} $\rightarrow$ \texttt{risk\_manager}, com aresta condicional ao \texttt{portfolio\_manager} apenas quando o veredito é \texttt{APROVADO}.}
\label{fig:grafo-langgraph}
{\footnotesize Fonte: elaborado pelos autores, a partir de \path{src/agents/graph.py}.}
\end{figure}

Os papéis são três. O agente de análise técnica (\texttt{technical\_analyst.py}) interpreta os indicadores quantitativos e emite uma recomendação inicial. O agente de gestão de risco (\texttt{risk\_manager.py}), amparado por um contrato de risco (\texttt{risk\_contract.py}), avalia a exposição e pode restringir a decisão. O agente de gestão de portfólio (\texttt{portfolio\_manager.py}) consolida as análises e define a decisão final. Na configuração científica, a direção da decisão é determinada pelo comitê e a exposição é atribuída por peso-alvo determinístico, sem dimensionamento dinâmico de posição pelo gestor de portfólio; o modo legado de dimensionamento por confiança permanece no código apenas como alternativa experimental.

O estado compartilhado do grafo é definido em \texttt{state.py}; as \textit{features} e sua semântica, em \texttt{features.py} e \texttt{feature\_semantics.py}; a evidência técnica apresentada aos agentes, em \texttt{technical\_evidence.py}; e os \textit{prompts} de análise técnica, em \texttt{technical\_prompt\_v4.py} e \texttt{technical\_prompt\_v5.py}. A comunicação com o modelo é encapsulada em \texttt{llm\_client.py}, com rastreamento em \texttt{llm\_trace.py}, e cada participante do comitê é representado em \texttt{participant.py}. O método de decisão é o \textit{self-consistency}, em que múltiplos analistas votam e a decisão consolidada exige um limiar de consenso com voto obrigatório de todos os participantes; a configuração congelada é apresentada no Capítulo~\ref{cha:metodologia}.

O grafo (\texttt{build\_graph}, em \path{src/agents/graph.py}) é utilizado pelos participantes da arena experimental e pelos executores de ativo único (\texttt{scripts/run\_agent\_backtest.py}, \texttt{scripts/run\_agent\_daily.py}). O participante recusa universos com mais de um ativo, e seu estado mantém posição escalar, caixa e patrimônio de um único papel. Por essa razão, a carteira multiativos não instancia o grafo: reutiliza o nó do comitê técnico e os esquemas de resposta de \texttt{src/agents}, mas substitui o gestor de risco e o gestor de portfólio por papéis próprios.

\section{Camada de experimentos e proveniência}
\label{sec:arq-experimentos}

A execução e o controle dos experimentos de ativo único concentram-se em \texttt{src/experiments}. O módulo \texttt{spec.py} define a especificação do experimento; \texttt{runner.py} conduz a execução; \texttt{hardening.py}, \texttt{cal\_b.py}, \texttt{stress.py} e \texttt{phases.py} implementam as fases do desenho experimental; \texttt{anchors.py} define as âncoras temporais; \texttt{treatment.py} e \texttt{context.py} organizam o tratamento e o contexto de cada execução; e os módulos \texttt{h2\_evaluation.py}, \texttt{h2\_evaluation\_offline.py}, \texttt{h2\_evaluation\_production.py} e \texttt{h2\_evaluation\_manifest.py} tratam da avaliação do sistema multiagente e da geração do manifesto de proveniência.

A reprodutibilidade é garantida por um conjunto de identificadores registrados a cada execução: o \textit{commit} do código, o \textit{hash} da especificação (\texttt{spec\_hash}), o identificador do \textit{snapshot} de dados, o manifesto de execução e o rastreio das chamadas ao modelo de linguagem. As fases e os identificadores de caso não fazem parte do \texttt{spec\_hash}, de modo a não vincular a especificação a uma janela temporal específica.

O caminho multiativos reutiliza dessa camada o diário durável de chamadas (\texttt{CallBank} e \texttt{DurableGeminiClient}, de \texttt{h2\_evaluation\_production.py}) e o bloqueio de rede (\texttt{no\_network}, de \texttt{h2\_evaluation\_offline.py}), com identidades próprias, sem alterar os \textit{gates} do protocolo de ativo único. Sua proveniência é registrada nos manifestos de campanha e nos pacotes descritos no Apêndice~\ref{ap:empacotamento}.

\chapter{Empacotamento e Publicação das Campanhas}
\label{ap:empacotamento}

Este apêndice detalha o mecanismo de identidade, artefatos, exportação e publicação das campanhas do Replay Contínuo, resumido na Seção~\ref{sec:arq-replay}.

\section{Executores, manifestos e \textit{lock}}

Cada geração possui um executor próprio: \texttt{scripts/continuous\_replay.py} (identidade \texttt{HFL-CONTINUOUS-EXPLORATORY-REPLAY-V1}), \texttt{scripts/continuous\_replay\_v2.py} (\texttt{-V2}) e \texttt{scripts/continuous\_replay\_v3.py} (\texttt{-V3}), cada um com raiz de armazenamento, livro, orçamento e pacote exportado próprios. As versões v2 e v3 conferem o \textit{hash} do \textit{bundle} de mercado antes do uso e reutilizam contabilidade, livro, planejador, transporte e orçamento por vinculação de módulo (\texttt{pilot\_paper.bind}). O \textit{bundle} da campanha (\texttt{inputs/market.json} e \texttt{inputs/manifest.json}) é gerado por \texttt{scripts/continuous\_replay\_data.py}.

No início de cada execução, o executor monta um manifesto com a identidade, o modo (inferência real ou teste \textit{offline}), as sessões, as estratégias, as políticas de referência, o teto de orçamento e os \textit{hashes} SHA-256 de todas as fontes. Um manifesto diferente na raiz recusa a execução (\texttt{campaign\_identity\_changed}). A v1 copia ainda os módulos de livro, proventos e regras para uma pasta imutável da campanha (\texttt{engine/}) e confere, por \textit{hash}, que as estratégias SMA/Bollinger e de portfólio importadas coincidem com essas cópias. Um \textit{lock} exclusivo impede dois executores simultâneos sobre a mesma raiz. As referências podem ser produzidas isoladamente (\texttt{-{}-references-only}). Na mínima variância, o otimizador SLSQP tem a função-objetivo dividida por um escalar positivo (a maior variância da diagonal), o que preserva o minimizador e as restrições; o eventual recurso aos pesos iguais também fica registrado no recibo de decisão. O aquecimento dos indicadores não cria patrimônio anterior ao início da campanha.

\section{Artefatos por pregão}

Para cada pregão, o executor grava na pasta da sessão, nesta ordem: \texttt{before.json} (estado das carteiras antes do pregão, após o crédito dos proventos pagos até a referência); \texttt{result.json} ou \texttt{failure.json} (resultado do comitê de cada ativo, com a entrada e o diário de chamadas); \texttt{context.json} e \texttt{proposal.json} (contexto e proposta compilada do gestor); \texttt{execution.json} (plano de execução); e \texttt{day.json} (fechamento do dia, com a marcação das seis carteiras). Cada comitê e cada gestor possuem um diário de chamadas (\texttt{CallBank}), exportado como \texttt{provider\_journal.jsonl}. A auditoria do orçamento (\texttt{audit\_budget}) reconstrói todas as tentativas, confere o \textit{hash} de cada envelope e recalcula o custo a partir dos \textit{tokens}.

\section{Exportadores}

Os exportadores (\texttt{scripts/export\_continuous\_replay.py}, \texttt{\_v2.py} e \texttt{\_v3.py}) leem os recibos já gravados, sem inferência nem alteração do livro, e produzem: \texttt{explorer.json} com as sessões, carteiras, pareceres e operações; \texttt{calls.json} com as chamadas arquivadas; tabelas CSV de patrimônio diário, operações, decisões, proventos e chamadas; uma página HTML autônoma com Chart.js embutido e um carregador das chamadas sob demanda (\texttt{calls-data.js}), gerados pelo próprio exportador nas versões v2 e v3 e por \texttt{scripts/build\_replay\_explorer.py} na v1; e um \texttt{manifest.json} selado com o \textit{hash} da exportação e o resultado da auditoria de orçamento. O conjunto é distribuído como pacote ZIP.

\section{Publicação e rotas}

A instalação no servidor é feita por \texttt{scripts/publish\_replay.py}, que aceita apenas o pacote cujo SHA-256 está fixado no código para cada versão e com a lista exata de membros esperada. Os arquivos são extraídos em área temporária, conferidos e promovidos por renomeação no mesmo volume; um destino existente só é aceito se seus arquivos coincidirem com o pacote, e dados diferentes nunca são sobrescritos. Cada versão ocupa um diretório próprio, preservando as anteriores.


A consulta é servida pelas rotas \path{/laboratory/replay/}, \path{/laboratory/replay/v2/} e \path{/laboratory/replay/v3/} de \texttt{dashboard/server.py}, somente com leitura: o servidor entrega apenas uma lista fechada de dez arquivos exportados, confina o caminho ao diretório da versão, comprime as respostas grandes sob demanda e não executa agentes nem operações. O explorador compara as curvas das seis carteiras e permite inspecionar, por carteira e pregão, posições, votos, pareceres de risco, decisões, operações, proventos e as chamadas arquivadas ao modelo. Os textos produzidos pelos modelos são apresentados como evidência preservada, não como fatos verificados.
